\documentclass{ctexart}
\usepackage{booktabs}
% 门控示例用的玩具论文，所有数字均为虚构。
\begin{document}

\section{数据}

样本包括416所学校的12,480名学生。第2节报告估计结果。

\section{结果}

班级越小，考试成绩越高。基准估计中班级规模的系数为$-$0.142（表~\ref{tab:main}第(1)列）。使用工具变量后系数为$-$0.215，在百分之一水平显著。班级人数少五人，对应的成绩高出1.08分（5 $\times$ 0.2153），约为平均成绩的1.8\%。\footnote{平均成绩为61.37分。}
% 旧版本在上句中写的是 $-$0.207

加入控制变量后，工具变量估计的系数为$-$0.188，在百分之五水平显著。基准估计中，班级人数少五人对应的成绩增益的95\%置信区间为[0.45, 1.70]分。

\begin{table}[h]
\centering
\caption{班级规模与考试成绩}\label{tab:main}
% generated from results/current/estimates.csv; do not edit by hand
\begin{tabular}{lccc}
\toprule
 & (1) & (2) & (3) \\
 & OLS & IV & IV \\
\midrule
Class size & $-$0.142$^{***}$ & $-$0.215$^{***}$ & $-$0.188$^{**}$ \\
 & (0.031) & (0.064) & (0.071) \\
Controls & No & No & Yes \\
First-stage $F$ & & 38.42 & 35.10 \\
Pupils & 12,480 & 12,480 & 12,480 \\
\bottomrule
\end{tabular}

\par\footnotesize 括号内为标准误。星号：*** $p<0.01$，** $p<0.05$，* $p<0.10$。
\end{table}

\end{document}
